\section{Caractères de rotation, couleur, saveur et virtualité}

Soit \(\jmath:\Gamma_{\rm TPK}\to\Gamma_{\rm TPK}\) l’involution qui inverse l’orientation de chaque arête et permute les fibres au moyen de \(J_c:E_c\to E_{\jmath c}\). Un caractère de rotation est une représentation \(\rho_{\rm rot}\) du sous-groupe engendré par \(\jmath\) sur les sections ; ses sous-espaces propres de signes \(+1\) et \(-1\) sont des objets algébriques distincts d’une représentation physique du groupe de Lorentz.

Soit \(\mathcal C\) le groupe abélien des caractères de canal du transport. Deux sections possèdent la même \emph{saveur arithmétique} lorsque leurs caractères
\[
 \chi_s:\mathcal C\longrightarrow\CC^\times
\]
coïncident. L’identification de ces classes à des saveurs de fermions exige une représentation de l’algèbre des observables physiques et n’est pas introduite par la nomenclature.

Enfin, soit
\[
 \operatorname{Ext}_N(s)=
 \{s'\in\Gamma(\mathcal B_{\gamma'}):
 \gamma'|_{\{0,\ldots,N\}}=\gamma,\ s'|_{\{0,\ldots,N\}}=s\}.
\]
La section est stable à partir de la profondeur \(N\) lorsque les restrictions \(\operatorname{Ext}_{N+1}(s)\to\operatorname{Ext}_N(s)\) admettent une branche compatible à tout niveau ultérieur. Elle est dite \emph{virtuelle jusqu’à la profondeur \(N\)} lorsque cette branche globale n’existe pas encore, bien que des extensions finies subsistent. Cette définition relève du système inverse et ne dépend pas d’une métaphore temporelle.

Sur le plan affine \(\FF_3^2\), la fonctionnelle linéaire
\[
 \kappa_3:\FF_3^2\longrightarrow\FF_3,
 \qquad \kappa_3(r,c)=r-c
\]
est surjective et son noyau est de dimension un. Ses trois fibres sont donc les classes affines du noyau et contiennent chacune exactement trois points. Le même équilibre persiste dans toute réunion disjointe de copies complètes de \(\FF_3^2\). Cette partition ternaire est une assertion algébrique autonome. La carte qutrit, l’orientation et le caractère de phase étant fixés, \cref{chap:algebra-color-qutrit} construit la représentation et son algèbre \(\mathfrak{su}(3)\) ; l’entrelaceur avec le transport TPK et l’identification chromodynamique physique demeurent une interface.
